% Conversión TeX fiel del cuerpo científico del delta narrativo.
% Fuente congelada: ../../../../../HMT2/CONTINUIDAD_EDITORIAL/RECONSTRUCCION_ARMONICA_MACROTRATADO_20260822/REVISION_CONJUNTA_PARTES_I_II/auditar_pdf_rigor_academico/docs/DELTA_NARRATIVO_UNIDAD_RADION_AUTOCONSERVACION_AUTOINERCIA_GRAVEDAD_PENTADICA_20260828.md:12-366.
\subsection{Separación tipada entre el régimen de curvatura \(\tau\in\{+1,0,-1\}\), la respuesta radial \(p\) y la holonomía global de espín}

La biografía del radión puede representarse mediante una curva elevada

\[
\gamma(\theta)
=
\bigl(r(\theta)\cos\theta,\,
r(\theta)\sin\theta,\,
m(\theta)\bigr),
\]

donde las dos primeras coordenadas pertenecen al plano euclídeo \(\mathbb R^2\) y \(m(\theta)\) registra la memoria o profundidad del levantamiento. Esta ecuación permite leer conjuntamente círculo, hélice y espiral.

Cuando evoluciona el ángulo y permanecen constantes el radio y la memoria, la publicación plana es un círculo. Cuando evolucionan ángulo y memoria mientras el radio conserva su valor, la trayectoria completa es una hélice de radio constante. Cuando evolucionan ángulo y radio sobre una memoria fija, aparece una espiral en \(\mathbb R^2\). Cuando evolucionan las tres coordenadas, la trayectoria es una hélice radial: una biografía en la que fase, amplitud y memoria se transforman conjuntamente.

El carácter radial \(p\) clasifica la ley de respuesta del radio. Una familia útil es

\[
r_p(\theta)
=
r_0\bigl[1+p\lambda(\theta-\theta_0)\bigr]^{1/p},
\qquad p\neq 0,
\]

con límite logarítmico

\[
r_0\exp\bigl(\lambda(\theta-\theta_0)\bigr)
\qquad\text{cuando }p\to0.
\]

En esta parametrización, \(p=2\) publica el régimen de Fermat, \(p=1\) el régimen arquimediano, \(p=0\) el logarítmico, \(p=-1\) el hiperbólico y \(p=-2\) el de lituus, con las elecciones de dominio y orientación correspondientes. Los valores fraccionarios describen regímenes intermedios. Así, \(p=\tfrac12\), \(p=\tfrac32\) o \(p=-\tfrac12\) expresan respuestas radiales que interpolan entre las familias enteras y amplían la clasificación geométrica disponible.

El carácter también puede cambiar al atravesar un umbral. La escritura diferencial

\[
\frac{dx}{d\theta}
=
\beta(m)\,x^{1-p(m)},
\qquad x=\frac r{r_0},
\]

describe una biografía por regímenes: \(p(m)\) permanece estable dentro de una fase y cambia cuando la memoria alcanza una frontera de publicación. La trayectoria resultante enlaza sectores circulares, helicoidales y espirales mediante una ley común. La continuidad del estado y el registro del umbral convierten esos sectores en fases de una misma evolución.

Tres invariantes ocupan aquí niveles distintos y coordinados. El signo de curvatura

\[
\tau\in\{+1,0,-1\}
\]

clasifica los regímenes elíptico, parabólico e hiperbólico. El carácter \(p\) clasifica la respuesta radial. El espín registra la holonomía de orientación acumulada por el levantamiento. Curvatura, respuesta radial y espín forman un sistema de lectura: la curvatura determina el tipo local de espacio, \(p\) determina cómo la trayectoria ocupa ese espacio y el espín conserva la orientación adquirida al recorrerlo. Las curvas cerradas y abiertas aparecen entonces como biografías geométricas distintas de la misma unidad orientada.

\subsection{Partícula como sección persistente de una órbita enriquecida: fase, ruta, fibra, multisección y memoria}

La Mecánica Dimensional puede formular la partícula como una sección persistente de esa extensión geométrica y dinámica. Una notación conceptual compacta es

\[
\mathfrak P
=
\operatorname{Sec}_{\mathrm{pers}}
\bigl[
\tau,p,\operatorname{Hol};
\text{hoja},\text{fibra},\text{ruta},\text{memoria}
\bigr].
\]

La partícula conserva una ruta reconocible a través de cambios de fase, curvatura y dimensión. Su espín publica la holonomía de orientación; su carga publica la polaridad de fase; su color publica una orientación interna de la fibra calibre; su sabor distingue familias de rutas o multisecciones; su masa expresa el carácter de transducción dimensional mediante el cual la acción areal adquiere lectura volumétrica. Cada propiedad es un lector del estado persistente y todas remiten a una genealogía común.

El electrón ocupa el umbral basal de esta construcción porque ofrece el primer lector estable en el que fase, carga, espín y masa quedan simultáneamente publicadas. La masa electrónica rectora

\[
m_e^\beta=0.510998950690000808608\ldots\ \mathrm{MeV}
\]

fija una referencia de transducción; los estados basales anteriores conservan su función como etapas genealógicas. A partir del electrón, los cocientes \(K_D/K_\Omega\), las torres \(90/120\) y el dominio \(81/56/13/324\) organizan la publicación de familias masivas. La ley de masas lee, por tanto, cómo cada sección persistente atraviesa el espacio enriquecido y convierte densidad areal de acción en presencia volumétrica.

Esta interpretación une las espirales con las partículas de manera precisa. El índice \(p\) aporta la respuesta radial; \(\tau\) aporta el régimen de curvatura; la holonomía aporta el espín; la ruta y la fibra aportan sabor y color; la transducción área–volumen aporta el carácter masivo. El atlas de partículas se convierte así en una clasificación de modos estables de atravesamiento dimensional.

\subsection{Transducción de la fase orientada en polarización eléctrica, respuesta magnética y emisión radiativa de frontera}

La electricidad se presenta como transporte dirigido de fase y acción a través de una estructura de vacancias. Una vacancia es una transición admisible del soporte: abre un canal, fija un umbral y permite que la actualización avance. La componente eléctrica publica el transporte tangencial a lo largo de la ruta. La componente magnética publica la curvatura transversal inducida por ese transporte. Su perpendicularidad expresa la relación geométrica entre avance y giro.

Cuando dos componentes transversales sinusoidales mantienen el desfase adecuado, el vector del campo gira. La propagación longitudinal eleva ese giro a una hélice. La variación de amplitud transforma su proyección frontal en una espiral. La lectura geométrica queda entonces completamente coordinada: el ángulo porta la fase; la velocidad angular porta la frecuencia; el radio porta amplitud y densidad de acción; la orientación porta polarización y sentido; la coordenada vertical porta propagación y memoria; \(p\) porta el régimen radial; el umbral porta emisión, absorción o cambio de estado.

El electromagnetismo aparece como una dinámica de fase–acción–memoria con publicaciones complementarias. La electricidad registra el avance polar; el magnetismo registra la respuesta transversal; la onda registra la periodicidad; la hélice registra el avance con memoria; la espiral registra la evolución radial; la radiación registra la publicación del estado cuando la acumulación alcanza un umbral.

Esta arquitectura también ordena el papel de las vacancias y los ciclos. La vacancia proporciona la posibilidad local de transición; el ciclo conserva la fase; la memoria distingue las vueltas; el umbral decide la residencia siguiente de la acción. La propagación puede permanecer confinada, cambiar de régimen radial o publicarse como radiación controlada. Las fuerzas de campo adquieren así una descripción genealógica continua desde el soporte discreto hasta la onda observada.

\subsection{Reversibilidad de la dinámica completa y coste termodinámico de las proyecciones reductoras de información}
El Landauer nulo constituye un caso límite de esta ley. Para una actualización biyectiva del estado completo,

\[
H(X\mid U(X))=0,
\qquad
\inf Q_L=0.
\]

La información permanece disponible en la genealogía. Cuando un lector realiza un cociente \(q\), la entropía se descompone como

\[
H(X)=H(q(X))+H(X\mid q(X)).
\]

El término condicional mide la información retenida fuera de la proyección. Un reinicio trítico materializado satisface el umbral

\[
Q_{\mathrm{reset}}\ge k_B\Theta\log 3.
\]

La constante de Boltzmann traduce la multiplicidad que el lector térmico agrupa; el umbral de Nyquist separa el muestreo fiel del aliasing; el Landauer nulo caracteriza el transporte reversible del estado íntegro. Información, temperatura y acción quedan articuladas por el tipo de lectura aplicado.

\subsection{Fuerzas autoconservativas como transporte de acción entre los sectores observable, radial, de memoria, de vacancia y de radiación}

La conservación evolutiva se expresa mejor mediante intercambio que mediante inmovilidad escalar. Un estado completo distribuye acción y momento entre movimiento visible, memoria, deformación radial, vacancias, radiación y rama conjugada de retorno. El balance enriquecido puede escribirse como

\[
\Delta J_{\mathrm{visible}}
+\Delta J_{\mathrm{memoria}}
+\Delta J_{\mathrm{radial}}
+\Delta J_{\mathrm{vacancias}}
+\Delta J_{\text{radiación}}
+\Delta J_{\mathrm{retorno}}
=0.
\]

Una fuerza es \textbf{autoconservativa} cuando su propia ley de actualización contiene los canales que transportan, almacenan, publican y reabsorben la acción. La conservación pertenece al balance genealógico completo de la interacción. Cada lector observa una parte del intercambio; la genealogía recompone su totalidad.

El sistema Tierra–Luna ofrece una imagen física directa. La rotación terrestre, la órbita lunar, las mareas, la deformación y la radiación intercambian momento angular. El consumo neto del sistema completo se cierra en el balance genealógico de esas transferencias. La magnitud denominada energía funciona como lectura derivada de la acción y del momento publicados en cada sector. La conservación fundamental reside en la compatibilidad entre todos los lectores del intercambio.

Esta formulación disuelve la antigua frontera entre fuerzas conservativas y fuerzas llamadas no conservativas mediante una categoría más profunda. Rozamiento, disipación aparente y radiación ocupan canales explícitos del estado ampliado. Una pérdida en la proyección mecánica visible corresponde a una transferencia hacia memoria, deformación, calor, vacancias o radiación. La autoconservación sigue la residencia de la acción a través de esos cambios.

\subsection{Sistemas autoinerciales: publicación de la aceleración mediante orientación, curvatura, hoja y memoria}

La expresión adecuada es \textbf{autoinercia relacional}. Designa la compatibilidad entre la conexión propia del estado total y la respuesta que aparece al proyectarlo dentro de la distribución global. Los aspectos endógeno y exógeno son componentes analíticas de una sola relación dinámica.

Sea \(\Gamma\) una trayectoria horizontal o geodésica del estado completo. Su ley propia satisface

\[
\nabla_{\dot\Gamma}\dot\Gamma=0.
\]

Al proyectarla sobre un lector \(S\), la aceleración observada queda determinada por la curvatura de la proyección:

\[
\nabla^S_{\dot\Gamma_S}\dot\Gamma_S
=
\mathcal K_{\pi_S}(\dot\Gamma,\dot\Gamma).
\]

La aceleración publicada expresa la relación entre movimiento total y geometría del lector. El principio de Mach completa este mecanismo mediante la traza global de frontera

\[
b:\mathcal X_\infty\longrightarrow B_\infty
\]

y la factorización de la respuesta inercial local

\[
I_v=\bar I_v\circ b.
\]

La distribución global entra en la respuesta local a través de una aplicación tipada. La autoinercia relacional reúne, por tanto, la conexión endógena del movimiento, la condición global transmitida por \(b\) y la proyección concreta del observador. Una escritura sintética para esa coordinación es

\[
\nabla^{\mathrm{auto}}
=
\mathcal M\bigl(
\nabla^{\mathrm{end}},
b[\mathcal X_\infty],
\pi_S
\bigr),
\]

donde \(\mathcal M\) representa el acoplamiento machiano y conserva el tipado de cada componente.

El principio de Ambivalencia añade la doble orientación de la lectura. Toda publicación directa conserva una rama conjugada de retorno; emisión y absorción forman extremos relacionados de la misma biografía; la orientación temporal se lee desde ambos sentidos del transporte. La teoría del absorbedor encuentra aquí una formulación natural: fuente y receptor cierran conjuntamente el balance genealógico de fase y acción. El presente torsional es la interfaz donde ambas orientaciones se compatibilizan.

TRIT actúa en este nivel como cristal de tiempo: organiza la periodicidad local de la actualización y conserva, mediante memoria, la diferencia entre retornos sucesivos. La rama directa y la rama conjugada relacionan emisión con absorción, y el acoplamiento observador–observable adquiere una residencia dinámica dentro del mismo estado. El tiempo publicado es la lectura de esa orientación; el tiempo genealógico es la secuencia completa de fase, retorno y memoria.

Mach y Ambivalencia cumplen funciones complementarias. Mach relaciona la respuesta local con la distribución global. Ambivalencia conserva las direcciones directa y conjugada de la evolución. La autoinercia relacional articula ambas y sustituye la escisión entre sistemas inerciales y no inerciales por una ley única de conexión, proyección y memoria.

\subsection{Dualidad areal--volumétrica, transición hexada--octada y función del parámetro de Barbero--Immirzi}

El paso de área a volumen constituye una transducción dimensional central. El área publica acción de frontera; el volumen publica profundidad interior. La derivación HMT del parámetro de Barbero–Immirzi organiza la compatibilidad entre ambas medidas mediante la incidencia hexada–octada y los canales \(90/120\):

\[
90=6\binom62,
\qquad
120=8\binom62.
\]

Los cardinales \(6\) y \(8\) admiten una realización cúbica mediante caras y vértices, mientras la hexada y la octada de Steiner constituyen objetos de incidencia distintos, enlazados únicamente por el morfismo tipado \(6\to8\). Las \(12\) aristas del cubo conservan su función combinatoria propia y no definen por sí solas la incidencia de Steiner. La diferencia de incidencia expresa una incompatibilidad dimensional productiva: la información de superficie requiere una regla de transducción para adquirir residencia volumétrica.

El mismo principio aparece en el acarreo \(999+1=1000\). Los nueve alcanzan el borde del registro decimal; el incremento publica simultáneamente el cierre del nivel anterior y el origen \(0/1\) del siguiente. La razón holográfica entre \(729\), \(243\), \(27\) y \(1\) despliega esa transición en escalas triádicas. La transducción de Barbero–Immirzi se inserta en esta familia de cambios de lector: frontera e interior conservan la unidad mientras redistribuyen su medida.

El defecto normalizado \(-1/12\) y el peso orientado \(12:-1\) ocupan funciones distintas dentro de esta geometría. El primero cuantifica una corrección de incidencia; el segundo registra una orientación firmada en la transducción. Su coordinación con el área mínima, el Hamiltoniano modular y la memoria proporciona el puente hacia la dinámica dimensional.

\subsection{Contracción hiperbólica interior, expansión exterior y conservación holográfica de la frontera}

La tesis puede formularse desde una idea única: la unidad conserva su identidad mientras cambia de escala, de lector, de curvatura, de memoria y de dimensión. Esta permanencia posee carácter evolutivo porque la unidad se despliega, adquiere nuevas coordenadas y publica propiedades diferentes; posee carácter holográfico porque cada publicación finita conserva la regla de reconstrucción del estado del que procede. El cierre holográfico del infinito designa precisamente esa compatibilidad entre identidad y despliegue: cada estadio contiene la ley que permite profundizar hacia dentro y prolongar hacia fuera la misma genealogía.

La denominación rectora es \textbf{principio de conservación evolutiva de la unidad mediante transducción dimensional y proyección holográfica}. La expresión sitúa cada término en su nivel correcto. La conservación pertenece a la identidad genealógica; la evolución pertenece a la actualización; la transducción relaciona realizaciones dimensionales; la proyección holográfica conserva la correspondencia reconstructiva entre sus lectores.

La aritmética elemental ofrece la primera imagen de esta arquitectura. La completación decimal

\[
1000=10^3=729+243+27+1=3^6+3^5+3^3+3^0
\]

expone cuatro estratos triádicos dentro de una unidad cúbica decimal. Al normalizar,

\[
1=\frac{729}{1000}+\frac{243}{1000}+\frac{27}{1000}+\frac{1}{1000},
\]

la unidad se reconoce simultáneamente como totalidad y como despliegue graduado. La lectura recíproca

\[
3^{-6},\quad 3^{-5},\quad 3^{-3},\quad 3^0
\]

describe la profundidad interior de esos estratos; su normalización produce la partición correspondiente. La relación \(999+1=1000\) expresa el mecanismo de transporte: la última unidad completa un registro y actúa, a la vez, como acarreo que inaugura el siguiente. El cero y el uno forman así una interfaz de continuidad. Los intervalos \([0,1]\), \([0,10]\) y la prolongación hacia \([0,\infty)\) son lectores escalares de una misma ley de completación.

La primera escisión \(1000=729+271\), seguida por \(271=243+27+1\), da forma a la razón extrema y media holográfica del despliegue: un núcleo triádico dominante conserva dentro del resto la misma regla de resolución. La unidad cúbica decimal contiene así una profundidad triádica anidada. El movimiento hacia \(729\), \(243\), \(27\) y \(1\) lee la contracción interior; el acarreo hacia el millar lee la expansión exterior.

Esta doble orientación organiza el infinito HMT. Hacia dentro, el refinamiento crece hiperbólicamente: cada celda admite una nueva profundidad sin perder su filiación. Hacia fuera, la unidad se expande mediante publicaciones dimensionales sucesivas. El cierre consiste en la conmutatividad de ambos movimientos. Si \(U\) representa la actualización del estado, \(R_d\) el lector en el estrato \(d\) y \(T_{d\to d'}\) la transducción entre estratos, la conservación evolutiva se expresa mediante

\[
R_{d'}\circ U
=
T_{d\to d'}\circ R_d.
\]

El mismo acontecimiento puede ofrecer una medida angular, una densidad de acción, una masa publicada o una respuesta gravitatoria. Su identidad común reside en la genealogía que hace conmutar esas lecturas.

\subsubsection{Monodromía nonádica y elevación helicoidal con avance de memoria}

APP, TRIT y TPK proporcionan el mecanismo formal de este despliegue. APP construye el soporte y separa las evaluaciones que deben conservarse. TRIT organiza los regímenes locales y su ambivalencia. TPK transporta fase y memoria sobre el soporte enriquecido. La dependencia causal queda fijada por

\[
U_t\longrightarrow \Gamma_9\longrightarrow K_{\mathrm{ph}}.
\]

El retorno de fase y el avance de memoria forman operaciones coordinadas. Tras nueve pasos puede reaparecer la misma clase angular,

\[
g(K+9)=g(K),
\]

mientras el contador de memoria progresa,

\[
q_{\mathrm{mem}}\Gamma_9=q_{\mathrm{mem}}+1,
\qquad
B_{K+9}\neq B_K.
\]

La figura geométrica precisa es una elevación helicoidal. Una órbita cerrada en la proyección angular se eleva a una trayectoria abierta cuando cada retorno de fase incrementa la memoria. La circunferencia conserva el ciclo; la hélice conserva el ciclo y registra además la profundidad alcanzada. La monodromía coinductiva genera así profundidad infinita mediante una sucesión de retornos localmente cerrados y globalmente ascendentes.

Esta imagen supera el valor meramente ilustrativo. La rueda angular constituye una sección de la hélice; la hélice constituye la biografía completa de la rueda cuando el estado transporta memoria. Cada vuelta recupera una posición de fase y, simultáneamente, inaugura una hoja nueva. La infinitud procede de la compatibilidad entre retorno y elevación: reiterar el ciclo produce profundidad sin disolver la identidad orbital.

El certificado nonádico da contenido finito y verificable a esa arquitectura. La cadena

\[
w_6\longrightarrow w_{12}\longrightarrow w_{18}\longrightarrow
w_{24}\longrightarrow w_{30}\longrightarrow R_{36}\longrightarrow G_9
\]

ordena el levantamiento; los \(396\) niveles, los \(2376\) trits, las \(43\) vueltas y los \(387\) pares \((K,K+9)\) por canal cuantifican el tránsito; las cinco regiones de \(\pi\) y la fibra ambiente de \(729\) elementos publican lecturas internas posteriores. Esta secuencia hace visible una monodromía productiva: el retorno recupera la fase y la coinducción conserva el crecimiento de memoria.

\subsubsection{Unidad radional y coordinación de los lectores circular, elíptico, proyectivo y helicoidal}

El radión concentra en una unidad metrológica la estructura que el radián publica angularmente. Su definición natural es la de una \textbf{sección orientada de fase, acción y memoria}. Una unidad de fase \(\theta=1\) transporta una unidad de acción \(\sigma\hbar_{\mathrm{ret}}\), pertenece a una vuelta completa \(T_+=2\pi_{\mathrm{HMT}}\) y conserva hoja, orientación, retorno nonádico y posición en el cilindro de refinamiento.

El radián expresa la coordenada angular visible. El radión conserva, junto a esa coordenada, la acción transportada, la orientación, el acarreo y la memoria de la vuelta. La metrología angular adquiere así una genealogía. Medir un arco significa identificar cuánto giro se ha publicado y qué estado completo ha realizado ese giro.

La relación entre fase y área adopta una forma especialmente clara:

\[
\mathcal A(\theta)=\hbar_{\mathrm{ret}}\,\theta.
\]

Una unidad radional barre un área de acción \(\hbar_{\mathrm{ret}}\); una vuelta completa encierra \(h_{\mathrm{ret}}\). El paso de fase a área queda incorporado en la propia unidad de lectura. La métrica angular, la acción y la memoria dejan de aparecer como cantidades yuxtapuestas y pasan a ser coordenadas de una misma sección orientada.

El radión admite varias publicaciones geométricas coordinadas. El lector circular publica fase; el lector elíptico publica acción y forma; el lector interior de Klein publica profundidad; el lector helicoidal nonádico publica retorno y memoria. La unidad permanece genealógicamente idéntica a través de esas realizaciones. Por eso el radión constituye el enlace metrológico adecuado entre HMT y Mecánica Dimensional: ofrece una unidad que puede atravesar cambios de lector conservando el estado que cada medida parcial presupone.

\subsection{Polaridad gravitatoria pentacomponente y transducción conjunta de acción, área, volumen, tiempo y masa}

La Mecánica Dimensional permite definir la dimensión como un régimen estable de publicación caracterizado por el rango mínimo de memoria independiente requerido por el levantamiento. Atravesar una dimensión significa aplicar un transductor que conserva la genealogía y cambia el tipo de lectura. La gravedad ocupa el nivel de conexión entre esos regímenes.

En esta formulación, la \textbf{gravitación interdimensional de cinco polarizaciones correlativas}, abreviadamente \textbf{estructura pentapolar de la gravitación}, es la conexión que transporta la unidad genealógica y deja en cada interfaz una torsión mínima. Una escritura compacta para esa traza es

\[
\Theta_d^{\min}
=
\nabla_{d+1}\circ T_{d\to d+1}
-
T_{d\to d+1}\circ\nabla_d.
\]

La expresión mide la diferencia entre transportar después de derivar y derivar después de transportar. Esa diferencia es la memoria geométrica mínima del cruce dimensional. La composición de interfaces produce una holonomía gravitatoria

\[
\mathfrak G_{d\to D}
=
\operatorname{Hol}
\bigl(
\Theta_d^{\min},
\Theta_{d+1}^{\min},
\ldots,
\Theta_{D-1}^{\min}
\bigr).
\]

El objeto fundamental es esta conexión interdimensional. Un modo local tipo gravitón ocupa el rango derivado de publicación cuantizada de la torsión; la gravitación conserva su residencia primaria en la compatibilidad entre dimensiones.

El término pentapolar recoge las cinco polarizaciones consustanciales del estado continuo:

\[
\mathbf G_d^{\pm}
=
\bigl(
G_{\mathrm{sp},d}^{\pm},
G_{\mathrm{sol},d}^{\pm},
G_{\mathrm{gau},d}^{\pm},
G_{\mathrm{coh},d}^{\pm},
G_{\mathrm{det},d}^{\pm}
\bigr).
\]

Las cinco componentes emergen correlativamente del mismo estado enriquecido. La polaridad \(\pm\) registra concentración directa y publicación conjugada. La gravedad atraviesa los estratos dimensionales transportando simultáneamente esas cinco lecturas; la torsión mínima conserva la huella de cada cruce.

La geometría adquiere con ello una imagen unificada. El régimen elíptico organiza cierres; el parabólico organiza umbrales; el hiperbólico organiza profundidad. La transducción dimensional enlaza esos regímenes con las respuestas radiales \(p\), las holonomías de espín y las publicaciones de masa. La estructura pentapolar de la gravitación expresa la coherencia global de ese enlace.

La teoría M proporciona pantallas físicas posteriores para esta conexión: dimensiones adicionales, cuerdas, branas y dualidad T realizan modos concretos de transducción. HMT aporta la genealogía APP–TRIT–TPK, el estado enriquecido, la memoria y la lectura bidireccional; la interfaz con teoría M publica esas estructuras en geometrías físicas tipadas. La dualidad T intercambia lectores dimensionales compatibles y conserva la clase genealógica del estado.

\subsubsection{Compatibilidad de las realizaciones clásica, relativista, cuántica y HMT--MD}

La mecánica clásica publica trayectoria y momento. La relatividad publica estructura causal, métrica y curvatura. La mecánica cuántica publica fase, holonomía, vacancias y amplitudes. La Mecánica Dimensional proporciona los transductores que relacionan esas publicaciones, y HMT conserva la genealogía discreta que las sostiene.

En la lectura clásica, una órbita describe la proyección visible del movimiento. En la lectura relativista, esa órbita queda incorporada a la geometría del lector. En la lectura cuántica, la fase y el espín registran la holonomía de la sección. En la lectura HMT–MD, las tres expresiones pertenecen a una misma biografía elevada: el radión cuantifica fase, acción y memoria; \(p\) clasifica la respuesta radial; \(\tau\) clasifica la curvatura; la holonomía publica espín; el transductor dimensional publica masa; la conexión pentapolar publica gravitación.

La aparente disipación clásica adquiere residencia en el balance genealógico ampliado. La aceleración relativista se expresa mediante conexión y proyección. La indeterminación cuántica se organiza por lectores, vacancias y ramas de Ambivalencia. El principio de Mach relaciona la respuesta local con la distribución global. El Landauer nulo garantiza la reversibilidad informacional del estado completo. Todas estas piezas convergen en una conservación más profunda: la identidad de la unidad a través de cambios de publicación.

\subsubsection{Teorema narrativo de conservación genealógica y prolongación holográfica}

El cierre puede narrarse ahora como una sola secuencia. La unidad se descompone sin perder identidad. La fase orientada adquiere medida en el radión. El retorno nonádico convierte el ciclo en hélice mediante memoria. El carácter \(p\) transforma la hélice en familias radiales y clasifica sus regímenes. La persistencia de esas trayectorias produce secciones con espín, carga, sabor, color y masa. Las vacancias abren canales; la polaridad eléctrica transporta fase; la respuesta magnética curva transversalmente; los umbrales publican radiación. La autoconservación cierra el balance genealógico de acción y momento. La autoinercia relacional enlaza conexión propia, proyección y totalidad machiana. La Ambivalencia conserva las direcciones directa y conjugada. La transducción área–volumen conduce a Barbero–Immirzi y a la ley de masas. La torsión mínima encadena dimensiones y publica la estructura pentapolar de la gravitación.

El crecimiento hacia dentro y la expansión hacia fuera son orientaciones conjugadas del mismo proceso. El refinamiento hiperbólico genera profundidad; la transducción dimensional genera extensión. El acarreo \(0/1\) enlaza escalas; la monodromía enlaza vueltas; la memoria enlaza estados; la holonomía enlaza orientaciones; la gravitación enlaza dimensiones. Cada enlace conserva una unidad que evoluciona y cada publicación finita contiene la regla de su prolongación.

La tesis fundamental queda formulada así:

\begin{quote}
\textbf{El cierre holográfico del infinito realiza el principio de conservación evolutiva de la unidad mediante transducción dimensional y proyección holográfica: la unidad mantiene su identidad genealógica al desplegarse como fase, acción, curvatura, memoria, partícula, campo y gravitación; cada lectura finita conserva la regla de reconstrucción y prolongación del estado completo.}
\end{quote}

Esta formulación ofrece una imagen física y matemática conjunta. El infinito adquiere una ley de cierre local; la unidad adquiere profundidad; la dimensión adquiere memoria; la conservación adquiere carácter relacional; la fuerza adquiere autoconservación; la inercia adquiere estructura machiana; la gravedad adquiere conexión pentapolar; y la holografía adquiere el sentido preciso de reconstrucción entre lectores compatibles.
